\begin{resumo}
  O agendamento de consultas e atendimentos de saúde no Brasil
  ocorre predominantemente por telefone ou presencialmente:
  pesquisa da Doctoralia e TuoTempo (2022) aponta que 91\% dos
  centros médicos utilizam o telefone como principal canal,
  sendo que uma em cada três ligações não é atendida. Essa
  ineficiência impõe dificuldades tanto aos pacientes, que
  dependem do horário comercial para marcar atendimentos,
  quanto aos profissionais e clínicas de saúde, que gerenciam
  suas agendas por meio de planilhas e aplicativos de mensagens.
  Diante desse cenário, o presente trabalho propõe o
  desenvolvimento da Medicano, uma plataforma web de agendamento
  de serviços de saúde que conecta pacientes a clínicas e
  profissionais de diferentes especialidades, incluindo medicina,
  psicologia, psiquiatria, odontologia e nutrição. A plataforma
  incorpora um chatbot de triagem inteligente baseado em Large
  Language Models que, por meio de uma conversa guiada, identifica
  as necessidades do usuário e o direciona para a especialidade
  mais adequada. A solução foi desenvolvida como um monorepo
  utilizando React e TypeScript no frontend, Node.js com NestJS
  no backend, MongoDB como banco de dados NoSQL e Vercel AI SDK
  para integração com modelos de linguagem. A infraestrutura é
  composta por frontend hospedado no Amazon S3 com distribuição
  via CloudFront e backend em instância Amazon EC2 com proxy
  reverso Nginx, com segredos gerenciados pelo AWS Secrets
  Manager. O desenvolvimento seguiu metodologia ágil com sprints
  quinzenais, resultando em um produto mínimo viável que
  contempla os fluxos de cadastro e autenticação, triagem
  inteligente por chatbot, busca de profissionais e clínicas
  por especialidade e localização, agendamento de consultas e
  gestão de agenda. O trabalho demonstra a viabilidade técnica
  da plataforma e estabelece uma base sólida para evolução
  contínua do produto.

  \vspace{\onelineskip}
  \noindent\textbf{Palavras-chave:}
  agendamento em saúde. triagem inteligente. large language
  models. marketplace de saúde. computação em nuvem.
\end{resumo}